\part{Morphology}

\chapter{Nouns}

\section{Formation}

\latscript{-CI} denotes those exercising a profession or habitual involvement with the object or concept expressed by the underived noun.

\begin{lingexample}
	\ex \latscript{balıqçı} `fisherman' \derivedfrom{} \latscript{balıq} `fish' \\
		\latscript{oquwcu} `reader' \derivedfrom{} \latscript{oquw} `reading' \\
		\latscript{temirci} `(black)smith' \derivedfrom{} \latscript{temir} `iron'
\end{lingexample}

\latscript{-DAş} denotes entities with a common attachment to the attribute, item or concept expressed by the underived noun.

\begin{lingexample}
	\ex \latscript{qarındaş} `sibling' \derivedfrom{} \latscript{qarın} `womb' \\
		\latscript{wätändeş} `fellow countryman' \derivedfrom{} \latscript{wätän} `homeland' \\
		\latscript{yoldaş} `fellow traveller, companion, comrade' \derivedfrom{} \latscript{yol} `road, path'
\end{lingexample}

\section{Inflection}

Nouns are inflected for number, possession and case.
When multiple such suffixes co-occur, they are added in the order: number-possession-case.

\begin{lingexample}
	\ex \gll üy-ler-imiz-den \\
		house-\Pl{}-\Poss{}.\Fpl{}-\Abl{} \\
		\glt `from our houses'
\end{lingexample}

\subsection{Number}

Singular nouns are unmarked, while plural nouns are formed by suffixing \latscript{-lAr} onto the singular stem.
Generally, only count nouns are pluralised.

\begin{lingexample}
	\ex \latscript{almalar} `apples' \derivedfrom{} \latscript{alma} `apple' \\
		\latscript{balalar} `children' \derivedfrom{} \latscript{bala} `child' \\
		\latscript{üyler} `houses' \derivedfrom{} \latscript{üy} `house'
\end{lingexample}

However, mass nouns can also be pluralised, with the plural suffix denoting non-specificity or discrete masses of the referent, rather than individual entities.

\begin{lingexample}
	\ex \latscript{ballar} `honeys, (jars/types/etc.\ of) honey' \derivedfrom{} \latscript{bal} `honey' \\
		\latscript{suwlar} `waters, (bodies/masses/etc.\ of) water' \derivedfrom{} \latscript{suw} `water'
\end{lingexample}

The plural suffix is never used if plurality is indicated by other means, such as by numerals or adverbs of quantity.

\begin{lingexample}
	\ex \latscript{beş oquwcu} `five students' \\
		\latscript{köp kişi} `many people'
\end{lingexample}

\subsection{Possession}

Possession is expressed by the suffixes in table~\ref{tab:possessive-suffixes}, which encode the person and number of the possessor.
The third person is the exception, which does not distinguish between a singular or plural possessor.

\begin{smalltable}[\label{tab:possessive-suffixes}]{Possessive suffixes}{c|c|c|c}
	& \textbf{\First{}} & \textbf{\Second{}} & \textbf{\Third{}} \\
	\hline
	\textbf{\Sg{}} & \latscript{-(I)m} & \latscript{-(I)ñ} & \multirow{2}{*}{\latscript{-(s)I}} \\
	\cline{1-3}
	\textbf{\Pl{}} & \latscript{-(I)mIz} & \latscript{-(I)ñIz} & \\
\end{smalltable}

If present, the possessor would be inflected in the genitive case.

\begin{lingexample}
	\ex \gll seniñ ayağ-ıñ \\
		\Ssg{}.\Gen{} foot-\Poss{}.\Ssg{} \\
		\glt `your foot'
\end{lingexample}

Because they act as determiners, the possessive suffixes confer definiteness onto the nominals to which they are attached.
Consequently, if said nominal is used as a direct object, it must be inflected in the accusative case.

\begin{lingexample}
	\ex \gll Öz qalqan-ı-n tut-ma-dı. \\
		\Refl{} shield-\Poss{}.\Third{}-\Acc{} hold-\Neg{}-\Pst{}.\Def{} \\
		\glt `He wasn\apostrophe{}t holding his (own) shield.'
\end{lingexample}

\subsection{Case}

\begin{smalltable}{Case suffixes}{c|c|c|c|c|c|c|c}
	\multicolumn{2}{|c|}{} & \textbf{\Nom{}} & \textbf{\Acc{}} & \textbf{\Gen{}} & \textbf{\Dat{}} & \textbf{\Loc{}} & \textbf{\Abl{}} \\
	\hline
	\multicolumn{2}{|c|}{\textbf{Normal}} & \multirow{3}{*}{\latscript{-\zeromorpheme{}}} & \multirow{2}{*}{\latscript{-nI}} & \multirow{3}{*}{\latscript{-nIñ}} & \latscript{-GA} & \latscript{-DA} & \latscript{-DAn} \\
	\cline{1-2} \cline{6-8}
	\multirow{2}{*}{\textbf{Possessive}} & \textbf{\First{}/\Second{}} & & & & \multirow{2}{*}{\latscript{-(n)A}} & \multirow{2}{*}{\latscript{-(n)dA}} & \multirow{2}{*}{\latscript{-(n)dAn}} \\
	\cline{2-2} \cline{4-4}
	& \textbf{\Third{}} & & \latscript{-n[I]} & & & & \\
\end{smalltable}

\chapter{Adjectives}

\section{Formation}

\latscript{-GI} forms relational adjectives from uninflected nominals (see example~\ref{ex:gi-relational}).
This suffix can also attach to nominals in the genitive (as \latscript{-KI}, see example~\ref{ex:gi-genitive}) or locative (see example~\ref{ex:gi-locative}), forming adjectives pertaining to possession or location.

\begin{exe}
	\ex \label{ex:gi-relational} \latscript{bugüngü} `today\apostrophe{}s, present-day' \derivedfrom{} \latscript{bugün} `today' \\
		\latscript{qışqı} `wintry, hibernal' \derivedfrom{} \latscript{qış} `winter'
	\ex \label{ex:gi-genitive} \latscript{meniñki} `mine' \derivedfrom{} \latscript{meniñ} `my' \\
		\latscript{quşnuñqu} `that which is a/the bird\apostrophe{}s' \derivedfrom{} \latscript{quşnuñ} `bird\apostrophe{}s'
	\ex \label{ex:gi-locative} \latscript{Astanadağı} `that which is in Astana' \derivedfrom{} \latscript{Astanada} `in Astana' \\
		\latscript{üyümdegi} `that which is in my house' \derivedfrom{} \latscript{üyümde} `in my house'
\end{exe}

If adjectives in the latter two groups are inflected for case without any other intervening morphology, the possessive set of case markers must be used.

\begin{exe}
	\ex \latscript{bizniñkinden} `from ours' \derivedfrom{} \latscript{bizniñki} `ours'
\end{exe}

Otherwise, adjectives formed from \latscript{-GI} can be used attributively.

\begin{lingexample}
	\ex \gll Bu awıl-da-ğı ağaç-lar uzun. \\
		\Prox{} village-\Loc{}-\Rel{} tree-\Pl{} tall \\
		`The trees in this village are tall.'
\end{lingexample}

\section{Comparison}

\chapter{Numerals}

\section{Cardinal numerals}

\section{Collective numerals}

\section{Ordinal numerals}

\section{Distributive numerals}

\chapter{Pronouns}

\section{Personal pronouns}

\chapter{Postpositions}

\chapter{Verbs}

\section{Formation}

\section{Stems}

Every verb has two stems: positive and negative.
Positive stems express the basic idea of the action and are formed by simply removing the infinitive (see example~\ref{ex:positive-stem-formation}).
This is also the verb\apostrophe{}s basic stem.
Negative stems express the non-performance of an action and are formed by suffixing \latscript{-mA} to the positive stem (see example~\ref{ex:negative-stem-formation}).

\begin{lingexample}
	\ex \label{ex:positive-stem-formation} \latscript{kör-} `see' \derivedfrom{} \latscript{körmek} `to see' \\
		\latscript{yaz-} `write' \derivedfrom{} \latscript{yazmaq} `to write'
	\ex \label{ex:negative-stem-formation} \latscript{alma-} `not take' \derivedfrom{} \latscript{al-} `take' \\
		\latscript{berme-} `not give' \derivedfrom{} \latscript{ber-} `give'
\end{lingexample}

\section{Non-finite forms}

\subsection{Participles}

\section{Person agreement}

\begin{smalltable}{Personal suffixes}{c|c|c|c|c|c|c}
	& \multicolumn{3}{c|}{\textbf{\Sg{}}} & \multicolumn{3}{c|}{\textbf{\Pl{}}} \\
	\cline{2-7}
	& \textbf{\First{}} & \textbf{\Second{}} & \textbf{\Third{}} & \textbf{\First{}} & \textbf{\Second{}} & \textbf{\Third{}} \\
	\hline
	\textbf{Pronominal} & \latscript{-mAn} & \latscript{-sAn} & \multirow{2}{*}{\latscript{-\zeromorpheme{}}} & \latscript{-mIz} & \latscript{-sIz} & \multirow{2}{*}{\latscript{[-lAr]}} \\
	\cline{1-3} \cline{5-6}
	\textbf{Possessive} & \latscript{-m} & \latscript{-ñ} & & \latscript{-K} & \latscript{-ñIz} & \\
\end{smalltable}

\section{Finite forms}

\subsection{Indicative}

\subsection{Conditional}

\chapter{Adverbs}
